\subsection{Serie de tiempo de los puntajes factoriales y la excedencia}\label{sec:serie-tiempo}

Los tres puntajes factoriales de \cref{sec:analisis-factorial} y \texttt{excede\_anio\_5}
(\cref{sec:obj:objetivo}) se caracterizan aquí en el tiempo: no se reabre su
origen ni su validez, ya establecidos, sino su comportamiento como series.

\subsubsection{De día-estación a área metropolitana}

Antes de aplicar cualquier descomposición de serie de tiempo hace falta una
serie sin huecos, y el panel día-estación no lo es. Solo 2 de las 15
estaciones (\texttt{CE}, \texttt{SE3}) superan 85\,\% de cobertura de
día-estación válido sobre el calendario 2021--2025, y la peor,
\texttt{NO3}, llega a 25\,\%; la cobertura global es 68.0\,\%. Filtrar a las estaciones de mejor
cobertura no resuelve el problema, lo esconde en las tres que sobreviven al
corte.

La alternativa es agregar las 15 estaciones a un valor diario de área
metropolitana, lo que solo es razonable si las estaciones no son
independientes entre sí el mismo día. La correlación transversal (Pearson,
estación contra estación, mismo día) va de una media de 0.50 en ventilación a
0.85 en forzamiento fotoquímico, con un mínimo apenas negativo únicamente en
ventilación. Es coherente con la física: la temperatura y la radiación solar
son fenómenos de escala regional, y el viento depende de la topografía local
del valle. Esta correlación es también la violación de independencia más
grande del conjunto de datos, relevante para cualquier modelo que trate cada
fila día-estación como una observación --- incluida la limitación de
autocorrelación ya declarada en \cref{sec:limitaciones-clasificacion}.

Con la cobertura insuficiente por estación y la correlación transversal alta,
se agrega a nivel área metropolitana promediando las estaciones que reportan
cada día. El promedio recupera 99.94\,\% de cobertura del calendario diario;
un único día, el 2021-02-16, se queda sin ninguna estación y se imputa por
interpolación lineal --- solo porque la descomposición MSTL exige una serie
sin huecos, no porque el hueco deje de existir.

\subsubsection{Descomposición MSTL}

Se ajusta MSTL con periodos semanal (7) y anual (365, no 365.25: MSTL exige
periodos enteros y el corrimiento acumulado en 4.5 años es de un día,
despreciable) a las cuatro series de área metropolitana. El ciclo semanal es
más fuerte en combustión primaria --- el patrón entre semana / fin de semana
del tráfico y la actividad industrial --- y más débil en la fracción de
estaciones que excede el umbral, que diluye ese patrón al depender de varios
factores a la vez.

\subsubsection{Autocorrelación residual}

La prueba de Ljung-Box en los rezagos 7, 14 y 30 rechaza con $p<0.001$ en las
cuatro series, en ambas versiones (cruda y residual de MSTL) --- el mismo
problema que \cref{sec:analisis-factorial} documentó para Bartlett: con $n=1{,}826$
observaciones diarias, el estadístico crece con $n$ y rechaza cualquier
autocorrelación no trivial, incluida la del residual donde la estacionalidad
ya se removió. La prueba no discrimina; la magnitud de la ACF sí
(\cref{fig:acf-excedencia}).

\begin{figure}[t]
  \centering
  \includegraphics[width=\linewidth]{acf_pacf_frac_excede_anio_5.pdf}
  \caption{ACF y PACF hasta el rezago 40 de la fracción que excede: serie
    cruda (arriba) contra residual de la descomposición MSTL (abajo), con
    banda $\pm 1.96/\sqrt{n}$.}
  \label{fig:acf-excedencia}
\end{figure}

Tras MSTL, la autocorrelación en el rezago 7 cae de 0.19--0.45 a valores
entre $-0.15$ y 0.05 en las cuatro series --- la estacionalidad semanal
prácticamente desaparece. La del rezago 1, en cambio, se mantiene entre 0.49 y
0.71 en el residual de las cuatro series. Esa persistencia de un día para
otro no es estacionalidad: es autocorrelación genuina de corto plazo, que
MSTL no está diseñado para remover. Cuantifica, con una serie continua y sin
huecos, la misma dependencia temporal que \cref{sec:limitaciones-clasificacion}
señaló como razón de que los errores estándar del modelo logístico sean
optimistas.

Dos consecuencias pasan al modelo autorregresivo de
\cref{sec:autorregresivo}: los rezagos deben calcularse sobre la serie de
área metropolitana reindexada a calendario diario regular, nunca sobre la
tabla día-estación, donde la fila anterior puede estar a una semana de
distancia; y la dependencia de un día para otro es genuina, no un artefacto
de estacionalidad, así que un modelo que la ignore deja estructura en los
residuales.
